\section{Optimization Problem: Soft-constrained optimal control}

\begin{frame}[more-space]{Optimization Problem: Soft-constrained optimal control}
  \dataheading{Trade computing capacity against cooling effort}
  {\footnotesize
  At measured inlet temperature, choose CPU power $P$, fan signal $F$ and
  thermal slack $s_T\geq0$. Proposed formulation:
  \[
    \begin{aligned}
      \max_{P,F,s_T}\quad & P-\lambda_F F-\rho s_T^2\\
      \text{subject to}\quad
      &\mu_T(P,F)+k\sigma_T(P,F)
          \leq T_{\mathrm{limit}}-\Delta T_{\mathrm{safety}}+s_T,\\
      &P_{\min}\leq P\leq P_{\max},\qquad F_{\min}\leq F\leq F_{\max},\\
      &|P-P_{\mathrm{before}}|\leq\Delta P_{\max},\qquad
        |F-F_{\mathrm{before}}|\leq\Delta F_{\max},\\
      &(P,T_{\mathrm{inlet}}^{\mathrm{meas}},F)\in\mathcal D_{\mathrm{supported}},
        \qquad s_T\geq0.
    \end{aligned}
  \]
  }
  \begin{columns}[T,onlytextwidth]
    \column{0.48\textwidth}
      \dataheading{Meaning of the objective}
      \footnotesize
      Reward CPU power as a computing-capacity proxy; penalize cooling effort
      and thermal-target violations.\par
      \smallskip
      $\lambda_F\geq0$ sets the cooling compromise; $\rho>0$ penalizes slack.
      Weight units/scales must match the objective.
    \column{0.49\textwidth}
      \dataheading{Why this is soft-constrained}
      \footnotesize
      Positive $s_T$ allows a penalized violation of the modeled thermal target.
      Actuator limits, step limits and domain restrictions remain hard.\par
      \smallskip
      Previous CPU/fan settings appear in step constraints, not as an initial thermal state of the PCE.
  \end{columns}
  \datasource{Proposed adaptation of the notes to the agreed soft-constraint objective.
    $F=P_{\mathrm{fan,raw}}$, unit unverified; $\lambda_F F$ is not a calibrated electrical-energy cost.
    Limits, weights, disturbance law and domain rule require specification; optimization is not implemented.}
\end{frame}

\begin{frame}[more-space]{Optimization Problem: Soft-constrained optimal control}
  \dataheading{Keep the optimizer inside conditions we can defend}
  \begin{columns}[T,onlytextwidth]
    \column{0.48\textwidth}
      \dataheading{A small, fast numerical problem}
      \footnotesize
      Evaluate the surrogate repeatedly for candidate CPU/fan settings.
      Use a small constrained solver; an explicit solution is optional only
      if the final formulation permits one.\par
      \smallskip
      The nonlinear surrogate and risk term may create local optima.
      Check candidate solutions, feasibility and sensitivity to starting points.
    \column{0.49\textwidth}
      \dataheading{Data and actuation limits}
      \footnotesize
      Observed ranges are an outer boundary, not proof that every combination
      inside the box was tested. Restrict $\mathcal D_{\mathrm{supported}}$
      using coverage and validation evidence.\par
      \smallskip
      Electrical fan power may not be directly commanded. A verified map from
      a supported fan setpoint to measured signal/airflow is needed for actuation.
  \end{columns}
  \medskip
  \begin{beamercolorbox}[sep=2mm,wd=\textwidth]{data warning}
    \dataheading{A surrogate optimum is a recommendation, not a hardware guarantee}
    \footnotesize
    An optimizer can exploit unrealistic predictions. Check physical trends and
    local prediction errors, then confirm the proposed setting with a validated
    full model or controlled experiment. Neither is currently available for this check.
    If evidence remains limited, present an offline demonstration rather than a validated controller.
  \end{beamercolorbox}
  \datasource{Project deployment/validation requirements from the notes.
    Coverage and fan-unit limitations: results/audit/report.md; results/inputs/report.md;
    data/raw/DATA\_DICTIONARY.csv. No optimal settings or solver results are claimed.}
\end{frame}
